\section{Related Works}
\label{sec:related}


% Kartik 11/13 start
\noindent\textbf{Preference Optimization for Med-LVLMs.} 
Preference optimization is widely adopted to align large vision–language models (LVLMs) with human or model preferences~\cite{Christiano2017RLHF,Schulman2017PPO,ouyang2022training,Ethayarajh2024KTO,liustatistical,hong2024reference,meng2024simpo,wu2025selfplay,liu2025survey}. More details are discussed in Appendix \ref{supp:related}. Extending such techniques to the medical domain introduces additional challenges, as factual grounding and clinical relevance are essential. Methods such as \cite{banerjee2024direct, sun2024stllava,hein2024preference} generate dispreferred responses via GPT-4 or the Med-LVLM for preference optimization.
Medical alignment emphasizes clinical grounding and lesion-aware signals.
MMedPO \cite{Zhu2025MMedPO} proposes clinically aware weighting to preference samples to highlight key medical cues, while MedGR$^{2}$ \cite{Zhi2025MedGR2} uses generative reward learning for robust medical reasoning.
Despite these advances, transferring general-domain pipelines still leaves gaps: models over-rely on textual priors and spurious background associations under imbalanced modality interactions \cite{Royer2024MultiMedEval,Cui2023GPT4VHallucination}.
MMedPO attempts to address these issues by weighting samples based on clinical relevance, but relies on computationally heavy multi-agent scoring and does not preserve useful lesion–background priors.
Guided by these observations, our work introduces a causally aware preference construction to probe lesion-mediated visual evidence while preserving clinically meaningful multimodal cues.

\noindent\textbf{Causal Reference.}
Causal inference \cite{pearl2009causality} uncovers cause–effect mechanisms through intervention and counterfactual reasoning.
Interventions modify distributions for unbiased estimates, while counterfactuals imagine hypothetical outcomes to for alternative conditions.
Recent learning-based approaches incorporate causal reasoning into deep models to reduce spurious shortcuts and improve generalization across vision tasks \cite{shao2024knowledge,li2024multimodal,zhou2025learning}, including anomaly detection \cite{liu2025crcl}, visual recognition \cite{liu2022contextual}, recommendation \cite{wang2022causal,wei2022causal}, and scene graph generation \cite{liu2025causal}.
In multimodal models, causal modeling is explored to improve robustness and interpretability.
Wang \textit{et al.} \cite{wang2025beyond} propose causal reward modeling to prevent reward hacking, but their framework is limited to unimodal text and does not extend to medical multimodal reasoning.
Zhou \textit{et al.} \cite{xu2025structure} integrate structural causal models into medical VQA to interpret visual–text relations, yet their front-door adjustment requires heavy sampling and fails to preserve lesion–background dependencies, rendering it suboptimal for multimodal alignment.
Building on these insights, we extend causal counterfactual reference to preference optimization. 
By constructing factual and counterfactual pairs, our framework suppresses harmful direct effects from biased supervision while preserving the valuable indirect effect that capture clinically meaningful lesion–background dependencies.

% Kartik 11/13 end



% OLD 
\iffalse
\paragraph{Preference Optimization for LMs/LVLMs.}
Preference optimization has become a cornerstone technique for aligning large vision–language models (LVLMs) with human or model preferences. Early work on reinforcement learning from human feedback (RLHF) \cite{Christiano2017RLHF} trains a reward model from pairwise human judgments and optimizes the policy through proximal policy optimization (PPO) \cite{Schulman2017PPO}. Direct Preference Optimization (DPO) \cite{rafailov2023direct} reformulates the alignment objective as a closed-form policy update derived from the Bradley–Terry model \cite{Bradley1952BT}, offering improved efficiency and stability. Several variants further refine the optimization objective or reward calibration strategy, such as Slic-DPO \cite{Zhao2023SlicDPO}, KTO \cite{Ethayarajh2024KTO}, RSO \cite{Liu2023RSO}, and ORPO \cite{Hong2024ORPO}, which introduce margin-based, utility-based, or sampling-based mechanisms for more consistent preference alignment. Variants such as SimPO \cite{meng2024simpo} and SPPO \cite{wu2025selfplay} reduce reliance on explicit reward modeling by simplifying objectives or generating preferences through self-play. In multimodal contexts, mDPO \cite{wang2024mdpo} extends DPO with conditional visual inputs such as multi-view or multi-region context, illustrating how pair construction and conditioning choices affect alignment. Our approach remains within the direct optimization family but introduces causal-aware sample weighting to better exploit clinically relevant visual evidence.

\paragraph{Preference Optimization for Med-LVLMs.} Recently, preference fine-tuning has also been extended to medical LVLMs, where factual grounding and clinical relevance are crucial. Works such as \cite{Banerjee2024MedAlign,Sun2024STLLaVAMed,Hein2024MedVQA} adapt preference optimization by generating dispreferred responses using GPT-4 or the target Med-LVLM itself. Medical alignment emphasizes clinical grounding and lesion-aware signals. MMedPO \cite{Zhu2025MMedPO} proposes clinical-aware weighting of preference samples to emphasize medically significant cues, while MedGR$^{2}$ \cite{Zhi2025MedGR2} explores generative reward learning for robust medical reasoning. Despite these improvements, transferring pipelines from general domains often leaves notable gaps: models may over-rely on textual priors and background spurious associations under imbalanced modality interactions \cite{Royer2024MultiMedEval,Cui2023GPT4VHallucination,Wu2023MedKLIP}. Although MMedPO attempts to address these issues by weighting samples based on clinical relevance, it depends on multi-agent scoring, which is computationally expensive and fails to preserve useful lesion–background priors. Guided by these observations, our work focuses on causal-aware preference construction that probes lesion-mediated visual evidence while retaining clinically meaningful multimodal dependencies.

\paragraph{Causal Reference.}
Causal inference \cite{pearl2009causality} aims to uncover cause–effect mechanisms, typically through intervention and counterfactual reasoning. The former manipulates variable distributions to obtain unbiased estimates, while the latter imagines hypothetical outcomes to reason about what would happen under alternative conditions. Recent learning-based approaches have incorporated causal reasoning into deep models to eliminate spurious shortcuts and enhance generalization across various computer vision tasks \cite{shao2024knowledge,li2024multimodal,zhou2025learning}, including video anomaly detection \cite{liu2025crcl}, visual recognition \cite{liu2022contextual}, recommendation \cite{wang2022causal,wei2022causal}, and scene graph generation \cite{liu2025causal}. In the context of large multimodal models, causal modeling has also been explored for improving robustness and interpretability. Wang \textit{et al.} \cite{wang2025beyond} introduce causal reward modeling to prevent reward hacking in LLM alignment, but their framework is limited to unimodal text and cannot be extended to medical multimodal reasoning. Zhou \textit{et al.} \cite{xu2025structure} integrate structural causal models into medical VQA to interpret visual–text relations, yet their front-door adjustment requires extensive sampling and fails to preserve informative lesion–background dependencies, making it sub-optimal for multimodal alignment. Building on these insights, we extend causal counterfactual reference to preference optimization. By constructing factual and counterfactual outcome pairs, our framework suppresses harmful direct effects from biased supervision while preserving the valuable indirect effect that captures clinically meaningful background-lesion dependencies.
\fi